% ---------------------------------------------------------------
%  Section III -- System Architecture and Problem Formulation
% ---------------------------------------------------------------

The IRIDIUM platform is organized into five functional layers, each with a
defined interface boundary and a set of registered data-source adapters.
Table~\ref{tab:layers} enumerates all layers, their constituent components,
the technology stack at the time of publication, and their deployment
readiness.
Table~\ref{tab:pipeline} describes the state-assembly pipeline that
implements the merge-only provenance policy.
These two tables replace a block-diagram representation deliberately: by
enumerating every component together with its runtime role and readiness
label, the architecture description becomes auditable and implementation
proximate, consistent with the provenance-first design philosophy of
IRIDIUM.

\begin{table*}[t]
\caption{IRIDIUM five-layer platform architecture. Each component is listed
         with its technology stack, deployment readiness, and activation
         state in the current release. ``Arch.'' denotes an architected
         component not yet operationally deployed.}
\label{tab:layers}
\centering
\renewcommand{\arraystretch}{1.20}
\begin{tabular}{@{}llllcc@{}}
\toprule
\textbf{Layer} &
\textbf{Component} &
\textbf{Technology} &
\textbf{Function} &
\textbf{Readiness} &
\textbf{Active} \\
\midrule
\multirow{7}{*}{External Sources}
  & Road network (OSM)          & Geofabrik PBF / HTTP & Topology data       & live              & Yes \\
  & Weather (Open-Meteo)        & REST / JSON          & Atmospheric data    & live              & Yes \\
  & Traffic (TomTom Flow)       & Licensed REST        & Speed telemetry     & config.\ req.     & No  \\
  & Transit (Baku Metro)        & GTFS / Agreement     & Schedule data       & perm.\ req.       & No  \\
  & Transit (BakuBus)           & GTFS / Agreement     & Schedule data       & perm.\ req.       & No  \\
  & Telemetry (Traccar)         & Consent-based GPS    & Crowd-sourced speed & not impl.         & No  \\
  & Events (iTicket)            & ToS review           & Event calendar      & not impl.         & No  \\
\midrule
\multirow{2}{*}{Ingestion \& Validation}
  & Schema validator            & Pydantic v2 / Python & Type-safe parsing   & live              & Yes \\
  & Provenance recorder         & Custom Python        & Lineage tracking    & live              & Yes \\
\midrule
\multirow{2}{*}{Digital Twin Assembler}
  & Graph builder               & PostgreSQL~16/PostGIS~3.4 & Topology merge & live              & Yes \\
  & State aggregator            & Python / NetworkX   & Dynamic state merge & live (partial)    & Yes \\
\midrule
\multirow{5}{*}{Analytics}
  & Congestion forecaster       & ST-GNN / heuristic  & Speed prediction    & heuristic only    & Partial \\
  & Multi-objective router      & Dijkstra / OTP (plan.) & Path optimization & baseline only     & Partial \\
  & Equity scorer               & NumPy               & $\MES_d$ computation & live             & Yes \\
  & Anomaly detector            & SciPy / Python      & Severity scoring    & live              & Yes \\
  & FL aggregator               & Flower (planned)    & FedAvg gradient agg.& arch.\ only       & No  \\
\midrule
\multirow{2}{*}{REST API}
  & API server                  & FastAPI / uvicorn~0.29 & HTTP interface  & live              & Yes \\
  & Status propagator           & Custom Python       & $\Pi$ attachment    & live              & Yes \\
\bottomrule
\end{tabular}
\end{table*}

\begin{table*}[t]
\caption{Digital twin state-assembly pipeline stages. Each stage produces a
         partial snapshot merged into the global graph $G$.
         Data-status labels propagate from acquisition through all
         downstream stages.}
\label{tab:pipeline}
\centering
\renewcommand{\arraystretch}{1.20}
\begin{tabular}{@{}clll@{}}
\toprule
\textbf{Stage} & \textbf{Operation} & \textbf{Input} & \textbf{Output} \\
\midrule
1 & Fetch            & Source adapter  & Raw bytes or HTTP response \\
2 & Validate         & Raw payload     & Typed Pydantic model \\
3 & Normalize        & Pydantic model  & Internal snapshot object \\
4 & Attach provenance& Snapshot + meta & $\pi_s$ record \\
5 & Merge into $G$   & Partial graph   & Updated $G$ (if \texttt{live} or \texttt{recorded}) \\
6 & Aggregate status & $\{\DS_s\}$     & $\DS_{\mathrm{global}}$ via Eq.~(\ref{eq:global_status}) \\
7 & Emit snapshot    & $G,\,\Pi$       & JSON API response \\
\bottomrule
\end{tabular}
\end{table*}

\subsection{Mobility Graph}

Let $G = (V, E, W)$ be a directed, attributed, multi-modal graph.
The node set
\begin{equation}
  V = V_\mathrm{seg} \cup V_\mathrm{jct} \cup V_\mathrm{stop}
\end{equation}
partitions into road segments, junctions, and transit stops.
Each node $v \in V$ carries a feature vector
\begin{equation}
  \mathbf{x}_v = [\text{type},\; \text{lat},\; \text{lon},\;
                  \text{mode}]^\top \in \mathcal{X}_v,
\end{equation}
where $\text{mode} \in \{\text{road, bus, metro, walking, cycling}\}$.
Each directed edge $e = (u,v) \in E$ carries a static attribute vector
\begin{equation}
  \mathbf{a}_e = [\ell_e,\; \tau_e^0,\; c_e,\; \kappa_e]^\top,
\end{equation}
where $\ell_e$ is edge length (m), $\tau_e^0$ is free-flow travel time (s),
$c_e$ is a generalized monetary cost (AZN), and $\kappa_e$ is a carbon
proxy (g CO$_2$-eq).
The optional dynamic state vector
\begin{equation}
  \mathbf{d}_e(t) = [s_e(t),\; \rho_e(t),\; \iota_e(t)]^\top
\end{equation}
carries current speed $s_e(t)$ (m\,s$^{-1}$), occupancy $\rho_e(t) \in [0,1]$,
and incident flag $\iota_e(t) \in \{0,1\}$.
The graph is assembled exclusively from configured real sources.
When no network import has been executed, $V = E = \emptyset$ and the
assembler returns $\DS = \texttt{config\_required}$.

The adjacency matrix $\mathbf{A} \in \{0,1\}^{|V|\times|V|}$ has
$A_{uv} = 1$ if $(u,v) \in E$.
A distance-weighted variant
\begin{equation}
  \tilde{A}_{uv} = \exp\!\left(-\frac{\ell_{uv}^2}{\sigma_\ell^2}\right)
  \cdot \mathbf{1}[(u,v)\in E]
  \label{eq:adj_weighted}
\end{equation}
with bandwidth $\sigma_\ell$ is used by the ST-GNN layer to encode
proximity-based spatial influence.

\subsection{Forecasting Formulation}

Let $\mathbf{X}_t \in \R^{|V| \times F}$ denote the node feature matrix
at time $t$ over $F$ features (speed, occupancy, weather conditions).
Given a lookback window of length $T$, the multi-step forecasting target is
\begin{equation}
  \hat{\mathbf{Y}}_{t+1:t+H} = f_\theta\!\left(
    \mathbf{X}_{t-T+1:t},\; G\right),
  \label{eq:forecast}
\end{equation}
where $H$ is the forecast horizon and $f_\theta$ is the parametric model.
Standard regression metrics for the held-out evaluation are
\begin{align}
  \mathrm{MAE}  &= \frac{1}{nH}\sum_{i=1}^{n}\sum_{h=1}^{H}
                   \left|y_{i,h} - \hat{y}_{i,h}\right|,
                   \label{eq:mae} \\
  \mathrm{RMSE} &= \sqrt{\frac{1}{nH}\sum_{i=1}^{n}\sum_{h=1}^{H}
                   \left(y_{i,h} - \hat{y}_{i,h}\right)^2},
                   \label{eq:rmse} \\
  \mathrm{MAPE} &= \frac{100\%}{nH}\sum_{i=1}^{n}\sum_{h=1}^{H}
                   \left|\frac{y_{i,h} - \hat{y}_{i,h}}{y_{i,h}}\right|.
                   \label{eq:mape}
\end{align}
The graph convolution step in a single ST-GNN layer reads
\begin{equation}
  \mathbf{H}^{(\ell+1)} = \sigma\!\left(
    \tilde{\mathbf{D}}^{-\tfrac{1}{2}}\tilde{\mathbf{A}}
    \tilde{\mathbf{D}}^{-\tfrac{1}{2}}\,
    \mathbf{H}^{(\ell)}\,\mathbf{\Theta}^{(\ell)}\right),
  \label{eq:gcn}
\end{equation}
where $\tilde{\mathbf{A}} = \mathbf{A} + \mathbf{I}$,
$\tilde{\mathbf{D}}$ is its degree matrix,
$\mathbf{\Theta}^{(\ell)}$ is the trainable weight matrix for layer $\ell$,
and $\sigma(\cdot)$ is a nonlinear activation (ReLU in the baseline
configuration).
The temporal dimension is handled by a gated recurrent unit (GRU) or
dilated causal convolution applied over the sequence
$\mathbf{H}^{(\ell)}_{t-T+1:t}$.

\subsection{Federated Aggregation (Architectural)}

Let there be $K$ client nodes, each holding a local dataset $\mathcal{D}_k$
of size $n_k = |\mathcal{D}_k|$.
The global FL objective is
\begin{equation}
  \min_{\theta}\; \sum_{k=1}^{K} p_k\, \mathcal{L}_k(\theta),
  \quad p_k = \frac{n_k}{\sum_{j=1}^{K} n_j},
  \label{eq:fl_obj}
\end{equation}
where
$\mathcal{L}_k(\theta) = \E_{(\mathbf{x},y)\sim\mathcal{D}_k}
[\ell(f_\theta(\mathbf{x}), y)]$.
The FedAvg global update after communication round $r$ is
\begin{equation}
  \theta^{(r+1)} = \sum_{k=1}^{K} p_k\, \theta_k^{(r)},
  \label{eq:fedavg}
\end{equation}
where $\theta_k^{(r)}$ is the parameter vector of client $k$ after $\tau$
local gradient steps.
Under the non-IID convergence result of Li et al.\
\cite{li2020fedavg_convergence}, the expected gradient norm satisfies
\begin{equation}
  \frac{1}{R}\sum_{r=0}^{R-1}
  \E\!\left[\norm{\nabla \mathcal{L}(\theta^{(r)})}^2\right]
  \leq
  \frac{2(\mathcal{L}(\theta^{(0)}) - \mathcal{L}^*)}
       {\eta\sqrt{RTK}}
  + \mathcal{O}\!\left(\frac{\eta\Gamma^2}{K}\right),
  \label{eq:fl_convergence}
\end{equation}
where $R$ is the total number of communication rounds, $\eta$ is the
learning rate, and $\Gamma^2$ quantifies the non-IID degree (gradient
divergence across clients).
The communication overhead per round scales as
$\mathcal{O}(K \cdot |\theta|)$ bits for uncompressed transfers, reducible
to $\mathcal{O}(K \cdot d \log b)$ for $d$-dimensional quantized updates
with $b$ bits per coordinate~\cite{kairouz2021advances}.

\subsection{Routing Objective}

A route $r = (e_1, e_2, \ldots, e_L)$ is an ordered sequence of edges
in $G$.
The scalar disutility to minimize is
\begin{equation}
  J(r) = \alpha\, T(r) + \beta\, C(r) + \gamma\, E(r) + \delta\, P(r),
  \label{eq:routing}
\end{equation}
where:
$T(r) = \sum_{e \in r} \tau_e(t)$ is realized travel time (s);
$C(r) = \sum_{e \in r} c_e$ is generalized monetary cost (AZN);
$E(r) = \sum_{e \in r} \kappa_e \cdot \ell_e$ is the emissions proxy
(g CO$_2$-eq); $P(r)$ is a mode-transfer or accessibility penalty;
$\alpha, \beta, \gamma, \delta \geq 0$ are user-configurable weights with
$\alpha + \beta + \gamma + \delta > 0$.
The feasibility constraint enforces a maximum journey time $T_{\max}$:
\begin{equation}
  \sum_{e \in r} \tau_e(t) \leq T_{\max},
  \quad r \in \mathcal{R}(o, d),
  \label{eq:routing_constraint}
\end{equation}
where $\mathcal{R}(o,d)$ is the set of all paths from origin $o$ to
destination $d$ in $G$.
When dynamic edge state $\mathbf{d}_e(t)$ is available from a configured
traffic provider, the realized travel time is
$\tau_e(t) = \ell_e / s_e(t)$; otherwise,
$\tau_e(t) = \tau_e^0$.
The computational complexity of Dijkstra's algorithm over $G$ is
$\mathcal{O}((|V|+|E|)\log|V|)$ with a Fibonacci heap priority queue,
which is tractable for Baku's road graph ($|V| \approx 10^5$,
$|E| \approx 2.5 \times 10^5$ from the OSM extract).

\subsection{Mobility Equity Score}

For district $d \in \mathcal{D}$ and accessibility indicator
$m \in \{1, \ldots, M\}$, let $x_{d,m}$ denote the raw indicator value
with empirical mean $\mu_m$ and standard deviation $\sigma_m$ over all
$|\mathcal{D}|$ districts.
The normalized indicator is
\begin{equation}
  z_{d,m} =
  \begin{cases}
    (x_{d,m} - \mu_m)\,/\,\sigma_m & \text{if } \sigma_m > 0, \\
    0 & \text{otherwise,}
  \end{cases}
  \label{eq:zscore_district}
\end{equation}
and the Mobility Equity Score for district $d$ is
\begin{equation}
  \MES_d = \sum_{m=1}^{M} w_m\, z_{d,m},
  \quad \text{s.t.}\;\sum_{m=1}^{M} w_m = 1,\; w_m \geq 0.
  \label{eq:mes}
\end{equation}
The citywide equity gap is
\begin{equation}
  \Delta_{\MES} = \max_{d \in \mathcal{D}} \MES_d
                - \min_{d \in \mathcal{D}} \MES_d,
  \label{eq:equity_gap}
\end{equation}
a scalar summary of spatial inequality in mobility access.
The variance of $\MES_d$ across districts,
\begin{equation}
  \mathrm{Var}(\MES) = \frac{1}{|\mathcal{D}|}
    \sum_{d \in \mathcal{D}}\!\left(\MES_d - \overline{\MES}\right)^2,
  \label{eq:mes_var}
\end{equation}
provides a complementary dispersion measure that is less sensitive to
outlier districts than $\Delta_{\MES}$.

\subsection{Anomaly Detection}

Let $x_t$ be the observed scalar signal (e.g., speed on a road segment)
at discrete time $t$.
Over a rolling window of length $w$, compute the rolling mean and standard
deviation:
\begin{align}
  \mu_{t,w}    &= \frac{1}{w}\sum_{i=0}^{w-1} x_{t-i},
                 \label{eq:roll_mean} \\
  \sigma_{t,w} &= \sqrt{\frac{1}{w-1}\sum_{i=0}^{w-1}
                   (x_{t-i}-\mu_{t,w})^2}.
                 \label{eq:roll_std}
\end{align}
The instantaneous z-score is
\begin{equation}
  z_t = \frac{x_t - \mu_{t,w}}{\sigma_{t,w} + \epsilon},
  \label{eq:zscore_anomaly}
\end{equation}
where $\epsilon > 0$ prevents division by zero when the signal is constant.
A composite anomaly severity score that fuses speed deviation, event
proximity, and adverse weather is
\begin{equation}
  S_t = \lambda_1 |z_t| + \lambda_2 e_t + \lambda_3 w_t,
  \quad \lambda_1 + \lambda_2 + \lambda_3 = 1,
  \label{eq:anomaly_severity}
\end{equation}
where $e_t \in [0,1]$ encodes the temporal proximity to a scheduled large
event and $w_t \in [0,1]$ encodes adverse weather intensity derived from
the Open-Meteo precipitation rate.
An anomaly is declared when $S_t > \tau_S$ for a configurable threshold
$\tau_S$.
Under the assumption that $z_t$ is approximately standard normal in the
absence of anomalies, the false-positive rate at threshold $\tau_S$ is
\begin{equation}
  \mathrm{FPR}(\tau_S) = 2\left(1 - \Phi(\tau_S / \lambda_1)\right),
  \label{eq:fpr}
\end{equation}
where $\Phi(\cdot)$ is the standard normal CDF.
For $\lambda_1 = 1$ (speed-only mode) and $\tau_S = 2.5$,
$\mathrm{FPR} \approx 1.24\%$ per time step.
